\documentclass[11pt]{article}
\usepackage[margin=1in]{geometry}
\usepackage{amsmath,amssymb}
\usepackage{hyperref}

\title{G8 Lane A Standard Eigenvalue Membership Route Audit}
\author{Private red-team memo}
\date{May 21, 2026}

\begin{document}
\maketitle

\section*{Executive Verdict}

The current Lean spine has reduced the nonzero-height side of Lane A to one
exact mathematical payload:
\[
  \forall z : \texttt{G8ActualXiNonzeroHeightCarrier},\quad
  \exists n : \mathbb{N},\quad
  V(z) = \lambda_n,
\]
where
\[
  V(z) =
  \texttt{g8ActualXiIotaTauCanonicalScaledOperatorSpectralValue}\ z
  = \iota_\tau^2\Bigl(s_z(1-s_z)-\frac14\Bigr)
\]
and
\[
  \lambda_n =
  \texttt{Tau.BookIII.Doors.lemniscate\_eigenvalue}\ n.
\]

Everything downstream from this exact membership theorem is already wired:
standard eigenvalue membership implies real-valuedness, real-valuedness of the
nonzero-height centered quadratic forces the critical axis, the zero-height
eta branch is theorem-backed, and the combined Lane A input gives actual
sigma-fixedness through the existing adapters.

The audit did not find an unconditional Book III theorem that already proves
the exact membership statement above.  Book III Chapters 23--25 prove and
motivate the operator and spectral-reality mechanism, but the actual zero to
eigenvalue membership step is routed through the determinant/spectral
correspondence bridge in Chapter 24.  That bridge is explicitly conditional in
the manuscript and is intentionally not used by the active Lane A route.

\section*{Current Lean Endpoint}

The smallest current proof-facing target is:
\[
\begin{aligned}
&\texttt{G8ActualXiCanonicalModeExtractionTarget} :
  \texttt{G8ActualXiNonzeroHeightCarrier} \to \mathbb{N},\\
&\texttt{G8ActualXiCanonicalModeEigenvalueAlignmentTarget}\ \texttt{modeOf} :
  \forall z,\,
  V(z) =
  (\texttt{Tau.BookIII.Doors.lemniscate\_eigenvalue}\ (\texttt{modeOf}\ z)
  : \mathbb{C}).
\end{aligned}
\]

These are packaged by
\[
  \texttt{G8ActualXiCanonicalStandardEigenvalueMembershipSource}.
\]
Once this package is supplied, the existing proof-checked adapters produce:
\[
\begin{gathered}
  \texttt{G8ActualXiIotaTauCanonicalStandardEigenvalueIndexSource},\\
  \texttt{G8BookIIILemniscateSpectralCoordinateRealizationSource},\\
  \texttt{G8CanonicalXiSpectralModeCorrespondenceSource},\\
  \texttt{G8ActualXiNonzeroHeightSpectralParameterReal}.
\end{gathered}
\]

\section*{Closed Lean Facts}

\paragraph{Zero-height side.}
The eta/paired-series path now supplies the zero-height axis guard.  In
particular, the positive-half-plane paired eta closure yields:
\[
  \texttt{g8ActualXiZeroHeightAxisGuardDischarge\_from\_pairedPositiveHalfPlane}.
\]

\paragraph{Centered quadratic algebra.}
For a nonzero-height carrier, real centered quadratic parameter forces
\[
  \operatorname{Re}(s_z)=\frac12.
\]
This is formalized by
\[
  \texttt{g8NonzeroHeightSpectralParameterReal\_forces\_re\_eq\_half}.
\]

\paragraph{Scale algebra.}
The certified \(\iota_\tau^2\) scale is nonzero, and the scaled-real to
unscaled-real adapter is already formalized.  Thus real-valuedness of
\[
  \iota_\tau^2\Bigl(s_z(1-s_z)-\frac14\Bigr)
\]
is sufficient to make the unscaled centered quadratic real.

\paragraph{Standard eigenvalue reality.}
Lean's standard lemniscate eigenvalue readout is definitionally \(n^2\), so a
value equal to a standard eigenvalue has zero imaginary coordinate after
complex readout:
\[
  \texttt{g8StandardLemniscateEigenvalueSpectrum\_real}.
\]

\paragraph{Coordinate plumbing.}
The finite coordinate system is now pure plumbing.  The self-stage coordinate
represents every standard mode:
\[
  \texttt{g8CanonicalXiSpectralCoordinateMode\_self\_stage}\ n = n,
\]
and the coordinate-realization target is equivalent to standard eigenvalue
membership.

\section*{Manuscript Anchors}

\paragraph{Chapter 23.}
The lemniscate operator is defined as \(H_L=-d^2/dx^2\) on the two-lobe metric
graph with continuity and Kirchhoff boundary conditions.  The manuscript states
self-adjointness of \(H_L\), and then obtains real, discrete spectrum from
standard compact metric graph theory.

\paragraph{Chapter 24.}
The spectral parameter is
\[
  \Lambda(s)=\iota_\tau^2\left(s(1-s)-\frac14\right),
\]
with reflection symmetry and critical-line mapping.  The chapter then states
the determinant representation and spectral correspondence as conditional:
\[
  \zeta_\tau(\rho)=0
  \quad\Longleftrightarrow\quad
  \Lambda(\rho)\in \operatorname{Spec}(H_L).
\]

\paragraph{Chapter 25.}
Assuming the spectral correspondence, the forcing argument expands
\[
  \rho(1-\rho)-\frac14
  =
  \left(\sigma(1-\sigma)+\gamma^2-\frac14\right)
  + i\gamma(1-2\sigma).
\]
Self-adjoint spectral reality makes the imaginary part vanish.  For
\(\gamma\ne0\), this forces \(\sigma=1/2\).

\section*{Line-by-Line Conditional Proof}

Assume the exact standard eigenvalue membership theorem:
\[
  \forall z,\ \exists n,\ V(z)=\lambda_n.
\]
Then:

\begin{enumerate}
\item Package the chosen mode \(n=\texttt{modeOf}\ z\) and equality
  \(V(z)=\lambda_n\) as
  \(\texttt{G8ActualXiCanonicalStandardEigenvalueMembershipSource}\).
\item Convert it to
  \(\texttt{G8ActualXiIotaTauCanonicalStandardEigenvalueIndexSource}\).
\item Use standard eigenvalue reality to conclude \(V(z).\operatorname{im}=0\).
\item Use the nonzero certified \(\iota_\tau^2\) scale to conclude
  \[
    \operatorname{Im}\left(s_z(1-s_z)-\frac14\right)=0.
  \]
\item Use the centered quadratic identity
  \[
    \operatorname{Im}\left(s_z(1-s_z)-\frac14\right)
    =
    \gamma_z(1-2\sigma_z).
  \]
\item Since \(z\) is in the nonzero-height subtype, \(\gamma_z\ne0\).
  Therefore \(1-2\sigma_z=0\), so \(\sigma_z=1/2\).
\item Thus no nonzero-height actual xi carrier is off-critical.
\item Combine with the already theorem-backed zero-height guard to obtain the
  height-split Lane A input.
\item The height-split input gives actual centered-address balance and actual
  sigma-fixedness.
\item With the separate accepted-tower realization-from-sigma-fixed input,
  the existing final live hinge and conditional handoff apply.
\end{enumerate}

\section*{What Is Still Missing}

The exact remaining mathematical theorem is:
\[
  \textbf{Actual standard eigenvalue membership:}\quad
  \forall z : \texttt{G8ActualXiNonzeroHeightCarrier},\
  \exists n : \mathbb{N},\
  V(z)=\lambda_n.
\]

Equivalently, one may prove the stronger source package:
\[
  \texttt{G8ActualXiCanonicalStandardEigenvalueMembershipSource}.
\]

This is stronger than merely proving \(V(z)\in\mathbb{R}\).  It asks for exact
membership in the standard discrete eigenvalue list.  The current finite
diagnostic checks do not prove it, and the manuscript route through Chapter 24
proves it only under the determinant/spectral-correspondence bridge.

\section*{Hidden-Gap Watchlist}

\begin{enumerate}
\item The Lean standard eigenvalue model uses
  \(\texttt{lemniscate\_eigenvalue}\ n=n^2\).  The manuscript metric-graph
  discussion also mentions secular families.  The next proof must verify that
  the chosen standard-mode list is the intended spectrum for this Lane A
  operator package, or weaken the source to an abstract spectrum membership
  predicate if exact \(n^2\) is too narrow.
\item Finite checks such as
  \(\texttt{self\_adjoint\_check}\),
  \(\texttt{spectral\_correspondence\_finite}\),
  \(\texttt{spectral\_resolution\_check}\), and
  \(\texttt{projector\_check}\)
  remain diagnostics.  They cannot construct the global mode \(n\).
\item Proving membership through O3, determinant transfer, accepted coverage,
  the final live hinge, full divisor transfer, or completion uniqueness would
  be circular for this Lane A source theorem.
\item If the intended route proves only real-valuedness of \(V(z)\), then the
  standard eigenvalue membership target is too strong.  The correct next Lean
  move would be to add an abstract self-adjoint spectral membership source
  rather than force equality to \(n^2\).
\end{enumerate}

\section*{Recommended Next Move}

Do not try to promote finite diagnostics into membership.  Instead choose one
of two honest next theorem targets:

\begin{enumerate}
\item Strong route:
  prove
  \[
    \texttt{G8ActualXiCanonicalModeEigenvalueAlignmentTarget}\ \texttt{modeOf}
  \]
  by constructing an exact standard mode for every actual nonzero-height
  carrier.
\item Weaker but possibly more faithful route:
  prove an abstract spectral membership theorem for the canonical scaled value
  in the self-adjoint lemniscate spectrum, then use spectral reality directly
  without requiring equality to the current \(n^2\) readout.
\end{enumerate}

The red-team recommendation is to attempt the weaker abstract membership route
first, because Book III Chapters 23--25 naturally support
\(\Lambda(s)\in\operatorname{Spec}(H_L)\) and spectral reality, while exact
equality to the finite \(n^2\) readout may encode additional discretization
choices that are not yet justified for every actual xi carrier.

\end{document}
